% GENERATED FILE. Edit canonical lessons, then run npm run generate:books.
% canonical-source-hash: fnv1a64:02098466d58e77c3
% canonical-lessons: TA-C292-vantal, TA-C292-malai-peytal, TA-R292-if-and-if-not

\chapter{If It Rains}
\label{ch:ta-if-it-rains}
\hlchaptermodality{292}

\begin{chapteropening}
\textbf{By the end of this chapter:} \emph{I can say what will happen if with vantāl and maḻai peytāl, say if after a noun with eṉṟāl, make an if less likely with oruvēḷai, and say even if (-ālum) and if not (-āviṭṭāl).}

Complete the last lesson of chapter 292: say what will happen if, even if and if not, with the if-form of a verb and eṉṟāl.
\end{chapteropening}

\section[\texorpdfstring{\ta{வந்தால்} (vantāl)}{வந்தால் (vantāl)}]{\ta{வந்தால்} (vantāl) — if}

\label{lesson:TA-C292-vantal}



\begin{quote}
Say \textquotedblleft{}I came\textquotedblright{}, \textquotedblleft{}I went\textquotedblright{} and \textquotedblleft{}I saw\textquotedblright{}. (\emph{vantēṉ}, \emph{pōṉēṉ}, \emph{pārttēṉ}.) Then \textquotedblleft{}in that case\textquotedblright{}. (\emph{appaṭiyāṉāl}.)
\end{quote}

\subsection*{You'll want to know: the if-form, -\ta{ஆல்}}
\begin{quote}
\textbf{\ta{நீங்கள்} \ta{நாளை} \ta{வந்தால்}, \ta{நாம்} \ta{சேர்ந்து} \ta{சாப்பிடலாம்}.}
\end{quote}

\begin{quote}
\emph{nīṅkaḷ nāḷai vantāl, nām cērntu cāppiṭalām.} — \textquotedblleft{}If you come tomorrow, we can eat together.\textquotedblright{}
\end{quote}

Build it from the \textquotedblleft{}I\textquotedblright{} past: take off \textbf{-\ta{ஏன்}}, put on \textbf{-\ta{ஆல்}}. \textbf{\ta{வந்தேன்}} $\to$ \textbf{\ta{வந்தால்}}; \textbf{\ta{போனேன்}} $\to$ \textbf{\ta{போனால்}}; \textbf{\ta{பார்த்தேன்}} $\to$ \textbf{\ta{பார்த்தால்}}; \textbf{\ta{இருந்தேன்}} $\to$ \textbf{\ta{இருந்தால்}}; \textbf{\ta{பேசினேன்}} $\to$ \textbf{\ta{பேசினால்}}. The past bead is inside it, but the if-form looks ahead: \textbf{\ta{வந்தால்}} is \textquotedblleft{}if (someone) comes\textquotedblright{}. It does not change for the person: \textbf{\ta{நான்} \ta{வந்தால்}}, \textbf{\ta{அவள்} \ta{வந்தால்}}.

You have said this shape inside \textbf{\ta{அப்படியானால்}}, \textquotedblleft{}in that case\textquotedblright{}: \textbf{\ta{அப்படி}}, \textquotedblleft{}so, like that\textquotedblright{}, with the if-form of \textbf{\ta{ஆகு}}, \textquotedblleft{}if it is so\textquotedblright{}. In speech the ending is clipped: \textbf{\ta{வந்தா}}.

\subsection*{You'll want to know: \ta{என்றால்}, and \ta{ஒருவேளை}}
After a noun, or a whole statement, \textquotedblleft{}if\textquotedblright{} is \textbf{\ta{என்றால்}}, the if-form of the verb inside \textbf{\ta{என்று}}: \textquotedblleft{}if (one) says\textquotedblright{}.

\begin{quote}
\textbf{\ta{நாளை} \ta{விடுமுறை} \ta{என்றால்}, \ta{நான்} \ta{ஊருக்குப்} \ta{போவேன்}.}
\end{quote}

\begin{quote}
\emph{nāḷai viṭumuṟai eṉṟāl, nāṉ ūrukkup pōvēṉ.} — \textquotedblleft{}If tomorrow is a holiday, I will go to my home town.\textquotedblright{}
\end{quote}

\textbf{\ta{ஒருவேளை}}, \textquotedblleft{}perhaps\textquotedblright{}, in front of an if makes it less likely: \textquotedblleft{}if by any chance\textquotedblright{}.

\begin{quote}
\textbf{\ta{ஒருவேளை} \ta{மழை} \ta{பெய்தால்}, \ta{குடை} \ta{எடுத்துப்} \ta{போங்கள்}.}
\end{quote}

\begin{quote}
\emph{oruvēḷai maḻai peytāl, kuṭai eṭuttup pōṅkaḷ.} — \textquotedblleft{}If by any chance it rains, take an umbrella with you.\textquotedblright{}
\end{quote}

\textbf{\ta{எடுத்துப்} \ta{போங்கள்}} is \textquotedblleft{}having taken, go\textquotedblright{}: the \textbf{-\ta{உ}} form, then the polite request.

\subsection*{Your turn}
Say these aloud:

\begin{itemize}
\raggedright
  \item \textquotedblleft{}If you come tomorrow, we can eat together\textquotedblright{} — \emph{nīṅkaḷ nāḷai vantāl, nām cērntu cāppiṭalām}
  \item \textquotedblleft{}If tomorrow is a holiday, I will go to my home town\textquotedblright{} — \emph{nāḷai viṭumuṟai eṉṟāl, nāṉ ūrukkup pōvēṉ}
  \item \textquotedblleft{}If by any chance it rains, take an umbrella with you\textquotedblright{} — \emph{oruvēḷai maḻai peytāl, kuṭai eṭuttup pōṅkaḷ}
\end{itemize}

\begin{tcolorbox}[breakable,colback=teal!4,colframe=teal!35!black,title={Before you move on}]
\textquotedblleft{}If (you) come\textquotedblright{}? (\textbf{\ta{வந்தால்}}.) \textquotedblleft{}If (you) go\textquotedblright{}? (\textbf{\ta{போனால்}}.) \textquotedblleft{}If it is a holiday\textquotedblright{}? (\textbf{\ta{விடுமுறை} \ta{என்றால்}}.) Which word makes the if less likely? (\textbf{\ta{ஒருவேளை}}.)
\end{tcolorbox}
\section[\texorpdfstring{\ta{மழை} \ta{பெய்தால்} (maḻai peytāl)}{மழை பெய்தால் (maḻai peytāl)}]{\ta{மழை} \ta{பெய்தால்} (maḻai peytāl) — if it rains, and if not}

\label{lesson:TA-C292-malai-peytal}



\begin{quote}
Say \textquotedblleft{}if (you) come\textquotedblright{}. (\emph{vantāl}.) Then \textquotedblleft{}even though\textquotedblright{} and \textquotedblleft{}otherwise\textquotedblright{}. (\emph{iruntālum}, \emph{illāviṭṭāl}.)
\end{quote}

\subsection*{You'll want to know: if this, then that}
\begin{quote}
\textbf{\ta{மழை} \ta{பெய்தால்} \ta{நான்} \ta{வீட்டில்} \ta{இருப்பேன்}.}
\end{quote}

\begin{quote}
\emph{maḻai peytāl nāṉ vīṭṭil iruppēṉ.} — \textquotedblleft{}If it rains, I will stay at home.\textquotedblright{}
\end{quote}

The if-clause comes first. The main clause after it says what follows: a future, \textbf{\ta{இருப்பேன்}} (\emph{iruppēṉ}, \textquotedblleft{}I will stay\textquotedblright{}); a request, \textbf{\ta{போங்கள்}}; or \textbf{-\ta{லாம்}}, \textbf{\ta{சாப்பிடலாம்}}.

Add \textbf{-\ta{உம்}} to the if-form and it becomes \textquotedblleft{}even if\textquotedblright{}:

\begin{quote}
\textbf{\ta{மழை} \ta{பெய்தாலும்} \ta{நான்} \ta{வருவேன்}.}
\end{quote}

\begin{quote}
\emph{maḻai peytālum nāṉ varuvēṉ.} — \textquotedblleft{}Even if it rains, I will come.\textquotedblright{}
\end{quote}

You know the shape from \textbf{\ta{இருந்தாலும்}}, \textquotedblleft{}even so, although\textquotedblright{}: the if-form of \textbf{\ta{இரு}} with \textbf{-\ta{உம்}}.

\subsection*{You'll want to know: if not, -\ta{ஆவிட்டால்}}
\begin{quote}
\textbf{\ta{நீங்கள்} \ta{சீக்கிரம்} \ta{வராவிட்டால்}, \ta{பேருந்து} \ta{போய்விடும்}.}
\end{quote}

\begin{quote}
\emph{nīṅkaḷ cīkkiram varāviṭṭāl, pēruntu pōyviṭum.} — \textquotedblleft{}If you don't come soon, the bus will leave.\textquotedblright{}
\end{quote}

Build \textquotedblleft{}if not\textquotedblright{} from the to-form: \textbf{\ta{வர}} $\to$ \textbf{\ta{வரா}-} with \textbf{\ta{விட்டால்}}: \textbf{\ta{வராவிட்டால்}}. So \textbf{\ta{போக}} $\to$ \textbf{\ta{போகாவிட்டால்}}, \textbf{\ta{சாப்பிட}} $\to$ \textbf{\ta{சாப்பிடாவிட்டால்}}. \textbf{\ta{போய்விடும்}} (\emph{pōyviṭum}, \textquotedblleft{}it will go off\textquotedblright{}).

\textbf{\ta{இல்லாவிட்டால்}}, \textquotedblleft{}otherwise\textquotedblright{}, is the same shape on \textbf{\ta{இல்லை}}: \textquotedblleft{}if it is not (so)\textquotedblright{}.

\subsection*{Your turn}
Say these aloud:

\begin{itemize}
\raggedright
  \item \textquotedblleft{}If it rains, I will stay at home\textquotedblright{} — \emph{maḻai peytāl nāṉ vīṭṭil iruppēṉ}
  \item \textquotedblleft{}Even if it rains, I will come\textquotedblright{} — \emph{maḻai peytālum nāṉ varuvēṉ}
  \item \textquotedblleft{}If you don't come soon, the bus will leave\textquotedblright{} — \emph{nīṅkaḷ cīkkiram varāviṭṭāl, pēruntu pōyviṭum}
  \item one plan for this weekend, with an if and an if not
\end{itemize}

\begin{tcolorbox}[breakable,colback=teal!4,colframe=teal!35!black,title={Before you move on}]
\textquotedblleft{}If it rains\textquotedblright{}? (\textbf{\ta{மழை} \ta{பெய்தால்}}.) \textquotedblleft{}Even if it rains\textquotedblright{}? (\textbf{\ta{மழை} \ta{பெய்தாலும்}}.) \textquotedblleft{}If you don't come\textquotedblright{}? (\textbf{\ta{நீங்கள்} \ta{வராவிட்டால்}}.) Which comes first, the if-clause or what follows? (\textbf{The if-clause.})
\end{tcolorbox}
\section[(review)]{Review — if, and if not}

\label{lesson:TA-R292-if-and-if-not}



\begin{quote}
Say \textquotedblleft{}if you come\textquotedblright{} and \textquotedblleft{}if you don't come\textquotedblright{}. Then read \textbf{\ta{ஈரம்}} aloud. (\emph{īram}.)
\end{quote}

\subsection*{Your turn}
Say these aloud:

\begin{itemize}
\raggedright
  \item If you come tomorrow, we can go to the sea.
  \item If it is a holiday, I will sleep until nine.
  \item If by any chance it rains, take an umbrella.
  \item Even if the bus is late, I will come.
  \item one if, one even if and one if not, about tomorrow
\end{itemize}

\begin{tcolorbox}[breakable,colback=teal!4,colframe=teal!35!black,title={Before you move on}]
\textquotedblleft{}If (you) come\textquotedblright{}? (\textbf{\ta{வந்தால்}}.) \textquotedblleft{}If it is a holiday\textquotedblright{}? (\textbf{\ta{விடுமுறை} \ta{என்றால்}}.) \textquotedblleft{}Even if it rains\textquotedblright{}? (\textbf{\ta{மழை} \ta{பெய்தாலும்}}.) \textquotedblleft{}If (you) don't go\textquotedblright{}? (\textbf{\ta{போகாவிட்டால்}}.)
\end{tcolorbox}
